\documentclass[12pt,conference,onecolumn]{IEEEtran}

\usepackage[hidelinks]{hyperref}

\title{Tennis biomechanics is fun}
\author{%
\IEEEauthorblockN{Miguel Arenas}\IEEEauthorblockA{Science \& Engineering\\Manalapan High School\\Englishtown, NJ\\\href{mailto:426marenas@frhsd.com}{426marenas@frhsd.com}}\and
\IEEEauthorblockN{Brady Gorelczenko}\IEEEauthorblockA{Science \& Engineering\\Manalapan High School\\Englishtown, NJ\\\href{mailto:426bgorelczenko@frhsd.com}{426bgorelczenko@frhsd.com}}\and
\IEEEauthorblockN{Timur Neyir}\IEEEauthorblockA{Science \& Engineering\\Manalapan High School\\Englishtown, NJ\\\href{mailto:426tneyir@frhsd.com}{426tneyir@frhsd.com}}}
\date{June 16, 2026}

\newcommand{\keywords}{biomechanics, sports, tennis, serve, kinematics}

\usepackage{hyperref}
\makeatletter
\AtBeginDocument{
\hypersetup{%
pdftitle={\@title},
pdfauthor={Miguel Arenas, Brady Gorelczenko, and Timur Neyir},
pdfkeywords={\keywords}}}
\makeatother

\begin{document}
\maketitle 

\begin{abstract}
The purpose of this research project is to improve our tennis serves by analyzing the shot velocity as we vary key components of the serve: leg bend, shoulder rotation, arm extension, and wrist pronation. The tennis serve we will be focusing on is a flat serve with a pinpoint stance and abbreviated backswing (see images). We aim to find the optimal versions of these components to maximize velocity. Some comparisons include: changing the angle behind the knees and vertical vs. horizontal-based shoulder rotation. To measure speed, we will either buy a radar gun or take a video of the ball's path and measure how much distance it covers between frames (similar to last year’s crossbow energy lab). \end{abstract}

\begin{IEEEkeywords}
\keywords
\end{IEEEkeywords}

\end{document}
